\section{UDP}
% Kurz und kanpp wechle Schicht was ist das.
Das UDP (User Datagram Protocol) ist ein verbindungsloses Protokoll der Transportschicht.
In diesem Versuch soll der verlustfreie und verlustbehaftete UDP Datenverkehr zwischen zwei Hosts untersucht werden.

\subsection{Durchführung}
Für die Durchführung wurde das Docker Netzwerk \ref{sec:aufbau} benutzt.
Es wurde zweimal der UDP Datenverkehr zwischen zwei Hosts untersucht.
Die Docker Umgebung wurde bei der Standardkonfiguration belassen. Einzig die Link-Parameter \ref{sec:link_einstellungen} wurden beim zweiten Durchlauf angepasst, um den Paketverlust zu simulieren:

\begin{itemize}[noitemsep]
    \item Bandwidth: 100 Mbit/s
    \item Delay: 0 ms
    \item Jitter: 0 ms
    \item Paketverlust: 30 \%
    \item Burst: 32kbit
\end{itemize}

In beiden Fällen wurden folgende Schritte durchgeführt:
\begin{enumerate}
    \item Starten der Docker Umgebung
    \item Einstellen der Link Eigenschaften \ref{sec:link_einstellungen}
    \item Starten des Netzwerkverkehr-Mitschnitts \ref{sec:pcap_docker_capture}
    \item Starten der UDP Verbindung zwischen den Hosts n1 (UDP Client) und n2 (UDP Server) \ref{sec:udp_docker_tool}
    \begin{itemize}[noitemsep]
        \item Server: n2
        \item Client: n1
        \item Count: 30
        \item Interval: 1 s
        \item size: 64 B
    \end{itemize}
    \item Beenden des Netzwerkverkehr-Mitschnitts \ref{sec:pcap_docker_capture}
    \item Schließen der Docker Umgebung
\end{enumerate}

Nach der Durchführung wurden die erzeugten PCAP Dateien mit Wireshark geöffnet und die UDP Pakete untersucht.

\subsection{Ergebnisse}
Bei dem Link ohne Paketverlust wurden alle 30 Pakete erfolgreich mit nahezu keiner Verzögerung übertragen.

\subsection{Diskussion}
% Welche Erkenntnisse konnten gewonnen werden mit bezug auf die Fragestellung?